\section{Methodology}\label{sec:methodology}
% Describe the systematic review protocol (e.g., PRISMA guidelines)

This systematic review was conducted in accordance with the Preferred Reporting Items for Systematic Reviews and Meta-Analyses (PRISMA) guidelines to ensure a rigorous, objective, and reproducible approach to identifying and selecting relevant literature.

\subsection{Search Strategy}\label{subsec:search-strategy}

A comprehensive literature search was performed using the databases Scopus, Web of Science, and IEEE Xplore to identify relevant studies. The search aimed to capture the intersection of space weather monitoring and artificial intelligence, specifically focusing on equatorial plasma bubbles. No date restriction was applied to the search, as the application of machine learning to EPB detection is a relatively recent development; the resulting temporal coverage of the included studies (2022--2025) reflects the natural emergence of the field rather than an imposed filter.

\subsection{Search String}\label{subsec:search-string}

The search string was constructed using Boolean operators to combine keywords related to the target atmospheric phenomena and the computational methodologies. The conceptual search string was (syntax was adapted to fit the specific formatting requirements of each database):

\begin{quote}
\small
\ttfamily
("plasma bubble*" OR "equatorial plasma bubble*" OR "EPB")
AND
("airglow" OR "all-sky imag*" OR "nightglow" OR "OI 630")
AND
("machine learning" OR "deep learning" OR "neural network*" OR "convolutional" OR "CNN" OR "classification" OR "detection")
\end{quote}

\subsection{Inclusion Criteria}\label{subsec:inclusion-criteria}

To determine the eligibility of the retrieved articles, the following inclusion criteria were established: 1. Peer-reviewed, original research articles. 2. Studies explicitly applying Machine Learning or Deep Learning algorithms. 3. Studies focused on the automated detection, classification, segmentation, or prediction of Equatorial Plasma Bubbles. 4. Studies utilizing relevant observational data (e.g., All-Sky Imager optical data, GNSS-TEC, or VHF radar). 5. Articles published in the English language.

\subsection{Exclusion Criteria}\label{subsec:exclusion-criteria}

Articles were excluded from the review if they met any of the following criteria: 1. Review papers, conference abstracts without full text, book chapters, or editorials. 2. Studies analyzing Equatorial Plasma Bubbles using solely traditional manual or statistical methods without Machine Learning components. 3. Studies applying Machine Learning to space weather phenomena unrelated to ionospheric plasma bubbles or disturbances. 4. Studies where the full text was inaccessible.

\subsection{Duplicate Removal and Data Management}\label{subsec:duplicate-removal-and-data-management}

Following the execution of the search strategy across the selected databases, all retrieved citations were exported to \textbf{Zotero}, a reference management software. Zotero was utilized to consolidate the records and automatically identify and remove duplicate publications, generating a list of unique records for the screening phase.

\subsection{Study Selection Process (Screening)}\label{subsec:study-selection-process-screening}

The screening process was conducted manually in two distinct stages by two independent reviewers to ensure a thorough and unbiased evaluation of the literature:

\begin{itemize}
    \item \textbf{Stage 1 (Title and Abstract Screening):} The titles and abstracts of all unique records were independently screened by both reviewers to assess their initial relevance to the research topic. Articles clearly falling outside the scope of the review were excluded at this stage.
    \item \textbf{Stage 2 (Full-Text Screening):} The articles that passed the initial screening were retrieved in full text and read in their entirety by both reviewers. A comprehensive independent evaluation of each manuscript was conducted against the predefined inclusion and exclusion criteria to confirm methodological relevance and definitively decide what to include or exclude for data extraction.
\end{itemize}

Both reviewers reached full agreement on the final set of included studies, requiring no conflict resolution.

\subsection{Number of Articles in Each Stage}\label{subsec:number-of-articles-in-each-stage}

The initial database search yielded an extensive pool of preliminary 37 records. After the automated duplicate removal process in Zotero, the remaining unique articles were subjected to the manual title and abstract screening. A subset of these articles advanced to the full-text screening phase. Upon rigorous manual full-text evaluation, articles failing to meet the strict inclusion criteria were excluded. Ultimately, exactly \textbf{9 articles} met all criteria and were included in the final qualitative synthesis. The complete study selection process is illustrated in Figure~\ref{fig:prisma-flow}.

\begin{figure}[ht]
\centering
\includegraphics[width=\textwidth]{figures/prisma-flow.drawio.pdf}
\caption{PRISMA flow diagram of the study selection process.}
\label{fig:prisma-flow}
\end{figure}